\section{Introduction}

The Web is entering a phase in which a growing share of interactions is performed by software on behalf of people. Language-model agents interleave natural-language reasoning with tool calls to search, browse, fill forms, send messages, and transact~\cite{yao2022react,schick2023toolformer}, and standardized tool interfaces such as the Model Context Protocol turn external services and data sources into first-class agent capabilities~\cite{hou2025model,zhang2025security}. Web-native benchmarks demonstrate agents operating realistic sites end to end~\cite{zhou2023webarena,deng2023mindweb}, tool-use suites exercise stateful, multi-step interactions~\cite{liu2023agentbench,lu2024toolsandbox}, API-based agents extend the reach of browsing~\cite{song2024beyond}, agentic search systems now return synthesized, citation-backed answers instead of link lists~\cite{gou2025mindweb}, and survey literature describes the resulting environment as an \emph{agentic Web} in which autonomous, goal-driven agents are primary actors~\cite{yang2025agentic,nie2026synergy,zhang2024large}. This shift reopens long-standing assumptions about who reads, writes, and pays for the Web, because the client is no longer a human reading a page but an agent executing operations---and every executed tool call can change external state.

The same shift changes the attack surface. Under indirect prompt injection, hostile instructions embedded in fetched content steer the model that processes them~\cite{greshake2023what,liu2023prompt,liu2024automatic}; when the steered model is an agent, the instruction is executed through tools with real consequences~\cite{zhan2024injecagent,debenedetti2024agentdojo,zhang2024agent}. Attacks and evaluations have followed this trajectory: environmental injection on generalist web agents~\cite{liao2024environmental}, prompt injection against multimodal web agents~\cite{wang2025webinject}, tool-layer attacks on protocol-mediated tool use~\cite{wang2026adaptools}, adaptive red-teaming of deployed web agents~\cite{syros2026muzzle}, and analyses showing that web agents are more vulnerable than standalone models precisely because each step multiplies exposure~\cite{chiang2025agents}. Attack success is commonly measured as an attack success rate over task runs~\cite{debenedetti2024agentdojo,zhan2024injecagent}, and dedicated safety benchmarks confirm that autonomous web agents routinely complete tasks in policy-violating ways~\cite{tur2025safearena,levy2024stwebagentbench}. Crucially for the Web, consequences are state changes---payments sent, records modified, private data exfiltrated---and recent work shows such consequences can be made covert, so the user never learns what their agent did~\cite{lee2026will}.

Current defenses address this threat at the wrong granularity, and their collective blind spots can be stated as two unanswered questions: whether cross-step execution dependencies are represented at all, and which tool operation actually produced the risk. The first question exposes segment-level methods. Content-level hardening and detection operate on text spans or attention patterns~\cite{chen2024struq,chen2024secalign,zhong2025attention}, guard models classify messages against safety taxonomies~\cite{inan2023llama,liu2024calibration}, and design-level systems constrain information flow between channels~\cite{debenedetti2025defeating,zhu2025melon}; all of them judge \emph{segments} in isolation, while the risk of an agent run is a property of \emph{compositions}: a benign query, a benign-looking tool output, and an otherwise reasonable plan can jointly produce an unauthorized transfer, and no per-segment verdict sees that propagation. The second question exposes run-level methods. Whole-trajectory judges score an entire execution~\cite{yuan2024rjudge,zhuge2024agentasajudge,mohammadi2026trajectoryjudge} and runtime monitors predict or veto upcoming steps~\cite{luo2025agrail,wang2025probguard,lin2026dreamguard,mou2026toolsafe,liu2026agentdog}; they answer \emph{whether} a run went bad, but not \emph{which} operation caused it, and outcome-only judges demonstrably miss process faults that do not change the final result~\cite{mohammadi2026trajectoryjudge}. An alert without a locus cannot become policy for the agent, cannot be audited by a site operator, and cannot be accumulated as knowledge; and since guardrails themselves are attackable and their verdicts are sensitive to formulation~\cite{zhou2026from}, evidence that is checkable and reusable has independent value.

We argue that the root cause of both gaps is a mismatch between the unit of risk and the unit of execution. Risk in tool-using agents propagates along real execution dependencies: external state enters through a tool result, is adopted by reasoning, and materializes in a tool call whose output becomes the next context. The natural boundary of this propagation is the \emph{action}: one tool call together with the state it consumed, the reasoning that justified it, and the state it produced. Reasoning has long been modeled as graphs in prompting frameworks~\cite{yao2023tree,besta2023graph}, and recent failure-attribution work shows that decisive errors can lie many steps away from their visible consequences, making dependency structure essential for locating them~\cite{deng2026defa}. Yet no existing representation binds risk evidence to complete tool operations: fixed-$k$-hop neighborhoods cut through operations at arbitrary points, keeping half of a reasoning chain while discarding the state that motivated it.

This paper makes the \aeu{} the anchor of agent safety analysis. We represent each execution as a directed \stg{} whose nodes are problems, thoughts, actions, and states, and whose typed and directed edges carry justification text. We construct, for every action, an inseparable \aeu{} and expand \csg{}s by adding whole predecessor units along real dependencies, so that every candidate is a complete, connected slice of an actual execution. A \csmil{} model encodes the full trajectory and all candidates with a shared edge-attention encoder, fuses local scores with the global score, and is trained with graph-level labels plus paired weak supervision derived from normal/anomalous executions of the same task, whose changed actions identify likely loci. At inference, a validated threshold triggers candidate ranking, and the top-ranked Action subgraph is compiled by a generator---which never sees ground-truth labels---into a structured \sexp{} that conditions the agent's subsequent decisions, closing the loop from detection to defense. Our contributions are:
\begin{itemize}
  \item a directed \stg{} representation for tool-using agents that jointly encodes tasks, reasoning, tool calls, and environment states, with the semantic type, direction, and justification of every edge;
  \item a \csg{} extraction and graph-level anomaly detection scheme in which \csg{}s are built on \aeu{}s---complete tool operations rather than fixed-hop neighborhoods---and a \csmil{} model that aligns full trajectories and candidates through a shared edge-attention encoder and, via learnable local--global aggregation and paired weak supervision, performs graph-level detection and Action-level localization jointly;
  \item a subgraph-level \sexp{} mechanism that compiles the top-ranked Action subgraph into a structured, node-grounded experience---generated without access to ground-truth labels---which conditions later agent decisions and closes the loop from detection to defense;
  \item an experimental design spanning detection, Action-level localization, and a controlled defense study whose conditions separate the effect of localized evidence from generic warning effects.
\end{itemize}
Section~\ref{sec:vision} develops why this formulation matters for the future Web: trajectory graphs as an auditable substrate and \sexp{}s as a portable safety layer. The rest of the paper formalizes the problem (\S\ref{sec:problem}), describes data and graph construction (\S\ref{sec:data}), details the method (\S\ref{sec:method}), specifies the evaluation design (\S\ref{sec:experiments}--\ref{sec:defense}), discusses implications for the agentic Web (\S\ref{sec:vision}), and concludes.
